\documentclass[a4paper,11pt]{article}
\usepackage[T1]{fontenc}
\usepackage{lmodern}
\usepackage[margin=0.85in]{geometry}
\usepackage{amsmath,amssymb,amsthm,mathtools}
\usepackage{xcolor,hyperref,enumitem}
\hypersetup{colorlinks=true,linkcolor=blue,urlcolor=blue}
\newtheorem{theorem}{Theorem}
\newtheorem{lemma}[theorem]{Lemma}
\newtheorem{corollary}[theorem]{Corollary}
\theoremstyle{remark}\newtheorem{remark}[theorem]{Remark}
\newcommand{\pp}[1]{\left(#1\right)_+}
\newcommand{\K}{\mathcal K}
\newcommand{\Om}{\Omega}
\setlength{\emergencystretch}{2em}
\title{A general two-bin certificate for Erd\H{o}s Problem 708}
\author{Research note: a proved lemma and the remaining arithmetic gap}
\date{24 November 2026; revised after the user's counterexample}
\begin{document}
\maketitle
\noindent\textbf{Revised status.} The general pointwise two-bin certificate remains proved. The unrestricted full-prefix transfer proposed in the first version is \textbf{false}: the user's construction, audited in Section~\ref{sec:user-counter}, gives a deficit of order $M^{10/17}$. The corrected own-full-tail estimate remains unproved. Neither $g(n)\le(2+o(1))n$ nor $g(n)\le2n$ is claimed.

\section{Setup and the new proved statement}
Let $M\ge1$ and let $I$ be any interval of $M$ consecutive positive integers. For finitely supported prices $0\le w_p\le1$, put
\[
 X(t)=\sum_p w_pv_p(t),\qquad h(t)=\prod_{w_p>0}p^{v_p(t)}.
\]
Thus $h(t)$ is the full active part, with every exponent retained. Primes larger than $M$ can be removed in the raw surplus problem. For $c=2+\delta>2$, write
\[
 A=\{a\le M:X(a)>c\},\quad n=|A|,\quad
 J_c=\sum_{a\le M}\pp{X(a)-c}-\sum_{b\in I}\pp{X(b)-1}.
\]
For a nonunit divisor $d$ supported on active primes, define
\begin{equation}\label{eq:root}
 r_w(d)=\max\left(0,\min\left\{\min_{p\mid d}w_p,
                   \frac{X(d)-1}{\Om(d)}\right\}\right),
 \qquad \Om(d)=\sum_pv_p(d).
\end{equation}
Set $r_w(1)=0$. For a scale $L\ge1$, let
\begin{equation}\label{eq:kernel}
 m_{p,L}(t)=\max_{\substack{d\mid h(t),\ p\mid d\\d^2\le Lp}}r_w(d),
 \qquad \K_L(t)=\sum_pv_p(t)m_{p,L}(t),
\end{equation}
where an empty maximum is zero.

\begin{theorem}[Uniform pointwise certificate]\label{thm:arith}
For every weight system as above,
\begin{align}
 \K_M(b)&\le\pp{X(b)-1} &&\text{for every positive integer }b,\label{eq:target}\\
 \pp{X(a)-2}&\le\K_{h(a)}(a)\le\K_a(a)\le\K_M(a)
                   &&\text{for every }a\le M.\label{eq:source}
\end{align}
There is no restriction on the number of active primes, their sizes or spacing, the number of occurrences, or the shift of $I$.
\end{theorem}

The source uses its \emph{own active size}. Replacing it by the ambient scale $M$ preserves the pointwise lower bound but can discard useful reserve; an exact example appears in Section~\ref{sec:failure}.

\section{Proof by two bins}
\begin{lemma}\label{lem:bins}
Let $V$ be a finite set, with weights $u_i\in[0,1]$ and sizes $\ell_i>0$. Put $U=\sum_i u_i$ and $L=\sum_i\ell_i$. For every nonempty $E\subseteq V$ define
\[
 r(E)=\max\left(0,\min\left\{\min_{j\in E}u_j,
                     \frac{\sum_{j\in E}u_j-1}{|E|}\right\}\right).
\]
In row $i$, permit all $E\ni i$ satisfying
\begin{equation}\label{eq:size}
 2\sum_{j\in E}\ell_j-\ell_i\le L,
\end{equation}
and let $m_i$ be the maximum $r(E)$ over these sets. Then
\begin{equation}\label{eq:bins}
 \pp{U-2}\le\sum_i m_i\le\pp{U-1}.
\end{equation}
The upper bound remains valid for arbitrary permitted families, chosen independently in the different rows.
\end{lemma}

\begin{proof}
Every $m_i$ is nonnegative and at most $u_i$. For the upper bound, there is nothing to prove if all maxima vanish. Otherwise choose a row and a permitted set $E$ attaining the largest row maximum, say $r>0$. For every $j\in E$, $m_j\le r$, regardless of whether $E$ is permitted in row $j$. Consequently
\[
 \sum_{j\in E}(u_j-m_j)
 \ge\sum_{j\in E}u_j-|E|r\ge1.
\]
All terms outside $E$ are nonnegative. This proves $U-\sum_i m_i\ge1$, and hence the upper bound using one common budget.

For the lower bound, set $v_i=u_i-m_i\in[0,1]$. Suppose, for a contradiction, that $\sum_i v_i>2$. Process the indices in nonincreasing order of $m_i$, placing items of size $v_i$ into either of two bins of capacity one. Use any rule when both bins can accept an item. At some first item $i$, placement must fail. If the current bin contents are $S$ and $T$, then
\begin{equation}\label{eq:failed}
 v_i+\sum_{j\in S}v_j>1,\qquad
 v_i+\sum_{j\in T}v_j>1.
\end{equation}
Also $m_j\ge m_i$ for every previously placed $j$. Let $E=\{i\}\cup S$ and $F=\{i\}\cup T$. These sets intersect only at $i$, and
\[
 [2\ell(E)-\ell_i]+[2\ell(F)-\ell_i]
   =2\bigl(\ell_i+\ell(S)+\ell(T)\bigr)\le2L.
\]
At least one of $E,F$, say $E$, satisfies \eqref{eq:size} in row $i$.

Remove from $E$ any $j\ne i$ for which $u_j=m_i$ and $v_j=0$. This preserves eligibility and the strict inequality $\sum_{j\in E}v_j>1$. The failed item has $v_i>0$, so it is not removed. Every remaining $u_j$ is strictly larger than $m_i$: indeed $u_j=v_j+m_j\ge v_j+m_i$, and the only possible equality case was removed. Moreover,
\[
 \sum_{j\in E}(u_j-m_i)
 =\sum_{j\in E}\bigl(v_j+m_j-m_i\bigr)
 \ge\sum_{j\in E}v_j>1.
\]
Both arguments in the minimum defining $r(E)$ are therefore strictly greater than $m_i$. This gives $r(E)>m_i$, a contradiction. Hence $\sum_i v_i\le2$, proving the lower bound.
\end{proof}

\paragraph{Arithmetic specialization.}
Represent every occurrence of a prime $p$ as a separate item of weight $w_p$ and size $\log p$. For a source $a$, their total size is $\log h(a)$. A subset representing a divisor $d$ and containing a designated occurrence of $p$ satisfies \eqref{eq:size} exactly when
\[
 2\log d-\log p\le\log h(a),\quad\text{that is, }d^2\le h(a)p.
\]
Any occurrence of the same prime can be designated, so its row maximum is exactly \eqref{eq:kernel}. The lower half of Lemma~\ref{lem:bins} gives \eqref{eq:source}. For a target $b$, use the arbitrary-family upper bound with families $d^2\le Mp$. This gives \eqref{eq:target}, even when $h(b)>M$. The case $h(t)=1$ has all quantities zero. This proves Theorem~\ref{thm:arith}.

\paragraph{Sharpness of the pointwise constant.}
Take $a=p^{2k}$, $k\ge2$, with the only positive price $w_p=1/k+\varepsilon\le1$. At scale $h(a)=a$, eligible powers $d=p^j$ have $2j\le2k+1$, hence $j\le k$. Their maximum price is $\varepsilon$. Thus
\[
 \K_{h(a)}(a)=2k\varepsilon=X(a)-2.
\]
The source inequality cannot gain a uniform additional positive amount per source solely from this pointwise construction.

\section{Exact common-budget interpretation}
If all subsets are permitted, the maxima have a simple closed form. For $U>1$, let $\theta\ge0$ be the unique solution of
\[
 \sum_i(u_i-\theta)_+=1.
\]
Then $m_i=\min(u_i,\theta)$ and $\sum_i m_i=U-1$. To prove this, first note that $r(E)\le\theta$ for every $E$: otherwise $\sum_{i\in E}(u_i-\theta)>1$. For $H=\{i:u_i>\theta\}$, the set $H$ attains $\theta$. If $u_i\le\theta$, adjoining $i$ to $H$ attains $u_i$. For $U\le1$, all root prices vanish.

The size restriction therefore allows a rigorously controlled loss of at most one further unit, from $U-1$ to $U-2$. This proves common target-budget feasibility and source coverage simultaneously. It does not count how often the permitted divisors occur in an interval.

For $c=2+\delta$, define
\begin{align*}
 R_A&=\sum_{a\in A}\bigl(\K_{h(a)}(a)-X(a)+2\bigr)\ge0,\\
 R_I&=\sum_{b\in I}\bigl((X(b)-1)_+-\K_M(b)\bigr)\ge0,\\
 \Gamma_c&=\sum_{a\in A}\K_{h(a)}(a)-\sum_{b\in I}\K_M(b).
\end{align*}
Then the exact identity is
\begin{equation}\label{eq:identity}
 \boxed{J_c=\Gamma_c-\delta n-R_A-R_I.}
\end{equation}
No prime-row balances have been transferred to a different source set or to a modified optimum.

\paragraph{The unproved step.}
The sufficient target is $\Gamma_c\le R_A+R_I+o_\delta(n)$ on genuine positive joint optimizers. It would contradict \eqref{eq:identity} for $n\to\infty$. The stronger assertion
\begin{equation}\label{eq:unproved}
 \sum_{a\le M}\K_{h(a)}(a)\le\sum_{b\in I}\K_M(b)
\end{equation}
was proposed as a stronger sufficient route in the first version. \textbf{It is false}, by Section~\ref{sec:user-counter}, even if an $O(\sqrt M)$ error is allowed. The own-full-tail assertion remains unproved. Equation~\eqref{eq:identity} records reserves; it is not itself an upper bound for $\Gamma_c$.

If \eqref{eq:unproved} held, Theorem~\ref{thm:arith} would give $\sum_{a\le M}(X(a)-2)_+\le\sum_I(X(b)-1)_+$. Boxed raw duality would imply $L_{\rm raw}(A,I)\le2|A|$. The raw finite-envelope theorem in Wang's manuscript would then give
\[
 g(n)\le2n+\pi(4n)-1=(2+o(1))n.
\]
This explains the motivation for the now-disproved route; it is not a remaining proof strategy. A valid uniform raw estimate can still use the existing transfer theorem without redoing rounding. The exact integer bound $g(n)\le2n$ remains separate.

\section{Two checks that prevent invalid transfer arguments}\label{sec:failure}
\subsection{Ambient-source transfer is false}
Take $M=18$, $w_2=w_3=1$, all other prices zero, and $I=\{9,\ldots,26\}$. Exact rational evaluation gives
\[
 \sum_{a\le18}\K_{18}(a)=\frac{17}{2}>8
 =\sum_{b=9}^{26}\K_{18}(b).
\]
In contrast,
\[
 \sum_{a\le18}\K_a(a)=\sum_{a\le18}\K_{h(a)}(a)=\frac{13}{2}<8.
\]
For $c=21/10$, the own full tail is $\{8,12,16,18\}$, with ambient kernel values $3/2,3/2,2,3/2$. Thus this example rejects an ambient-prefix shortcut, not the active-part transfer \eqref{eq:unproved} or its full-tail version.

\subsection{Independent row and price-level transfer is false}
Take $M=2N$, unit prices on all primes $p\le M$, and $2<c<3$. Then $A=\{a\le M:\Om(a)\ge3\}$. In row $p=2$ at price level $1/2$, every source $a=2x\in A$ has a permitted divisor $d=2q$: for the least prime factor $q$ of composite $x$, $q^2\le x$, so $d^2/2\le a=h(a)$.

The exact source row-level count, including all occurrences of two, is
\[
 S=\sum_{\substack{x\le N\\x\text{ composite}}}(1+v_2(x))
  =2N-s_2(N)-\pi(N)-2.
\]
The target event at $b=2x$ is that $x$ has a prime divisor $q\le\sqrt N$. For the fully explicit interval
\[
 N=2{,}000{,}000,\qquad
 I=\{2^{101}+1,\ldots,2^{101}+4{,}000{,}000\},
\]
an integer sieve gives
\[
 S=3{,}851{,}058,\qquad U=3{,}845{,}378,
 \qquad S-U=5{,}680>\sqrt M=2{,}000.
\]
The full calculation is included in the companion checker. More generally, averaging actual shifts and using the prime number and Mertens product asymptotics gives row deficits of order $N/\log N$, so no fixed square-root error works for independent row-level transfer.

This example is \emph{not} a positive joint optimizer; its tail is very dense. It does not refute \eqref{eq:unproved}. It shows why summing independently asserted row estimates would be invalid, even with the original full tail and the actual source scale.

\paragraph{A favorable arithmetic fact that does survive.}
For two eligible roots in the same row $p$,
\[
 d^2\le Mp,\quad e^2\le Mp
 \quad\Longrightarrow\quad
 \operatorname{lcm}(d,e)\le\frac{de}{p}\le M.
\]
Thus unweighted pair intersections are short. Higher intersections need not be short, and complete valuation layers involve $\operatorname{lcm}(p^j,d,e)$, which can exceed $M$. The whole-row balance does not authorize deleting those terms.

\section{An additional general restriction on positive surplus}
Let
\[
 \mu=\sum_{p\le M}\frac{w_p}{p-1},\qquad
 e=\mu-\frac1M\sum_{a\le M}X(a)\ge0.
\]
Prime-power counting gives
\[
 \mathbb E X^2\le\mu^2+\sum_p w_p^2\frac{p}{(p-1)^2}
 \le\mu^2+2\mu.
\]
Here and below expectation is over $1,\ldots,M$. For $x\ge0$ and $0\le q\le1$,
\[
 \min(x,2)\ge qx-\frac{q^2x^2}{8}.
\]
For $x\le2$ the right side is at most $qx\le x$; for $x\ge2$ the quadratic never exceeds two. If $\mu\ge2$, choose $q=4/(\mu+2)$. It follows that
\[
 \mathbb E\min(X,2)\ge\frac{2\mu-4e}{\mu+2}.
\]
Every prime-power has at least as many multiples in $I$ as in the prefix. Therefore, for $c\ge2$,
\[
 J_c\le J_2\le M-\sum_{a\le M}\min(X(a),2).
\]
We obtain the uniform necessary condition
\begin{equation}\label{eq:mean}
 \boxed{J_c>0\quad\Longrightarrow\quad\mu<2+4e
        =2+O(1/\log M).}
\end{equation}
The floor errors in $e$ are at most the number of prime powers up to $M$, divided by $M$; omitted geometric tails contribute $O(M^{-1/2})$. This proves the displayed uniform error. This mean restriction needs no optimality assumption, but does not prove the arithmetic transfer.

\section{Verification and scope}
The proof of Lemma~\ref{lem:bins}, not the experiments, establishes its uniformity. A companion exact-rational checker additionally verified 4,742 weighted-size configurations, including sharp families with arbitrarily many occurrences, and the $M=18$ arithmetic example. A separate integer checker recomputed the explicit row-level example above.

Exploratory searches in the first version tested both $\K_{h(a)}(a)$ and $\K_a(a)$ against actual CRT phase patterns and found no violation. The user's subsequent construction disproves \eqref{eq:unproved}; those finite searches missed its mechanism. They were not exhaustive global certificates and supplied no proof. They computed no new values of $g(n)$.

\noindent\textbf{What is proved:} the common-budget pointwise certificate, its two-bin proof and sharpness, the exact surplus identity, the stated counterexamples to stronger transfer shortcuts, and \eqref{eq:mean}.

\medskip
\noindent\textbf{What remains unproved:} the reserve-corrected own-full-tail estimate $\Gamma_c\le R_A+R_I+o_\delta(n)$ on genuine positive optimizers, or another sufficient arithmetic bound. The unrestricted full-prefix route is disproved. Consequently the general asymptotic and the exact $2n$ bound remain unresolved in this note.

\section{The user's counterexample to full-prefix transfer}\label{sec:user-counter}
\textbf{Attribution.} The construction in this section was supplied by the user. The audit makes the whole-prime multiplicities, CRT averaging, and small positive background prices explicit.

Let $p$ tend to infinity through primes and put
\[
 W=\lfloor p^{10/7}\rfloor,\quad M=pW,\quad
 Q=\{q\text{ prime}:W^{31/100}<q\le W^{8/25}\}.
\]
Temporarily use prices $w_p=1$, $w_q=1/100$ for $q\in Q$, and zero elsewhere. Denote the resulting kernel at a \emph{specified} scale $L$ by $\K_L^0$. The source below is $\K_a^0(a)$, not $\K_{h_0(a)}^0(a)$. Positive background prices will subsequently make $h(a)=a$.

\subsection{All nonzero main-price witnesses}
For large $p$,
\[
 p=W^{7/10}(1+o(1)),\quad W<p^2,\quad pq>W,
 \quad q^2<W,\quad(qr)^2>W\qquad(q,r\in Q).
\]
At target scale $M$, the only positive witnesses using main primes are $d=pq$ in row $p$, with price $1/200$. Indeed, a pure $Q$-witness has at most three occurrences: four would give $d^2>W^{2.48}$, whereas $Mq=W^{2.02+o(1)}$. Its load is then below one. The divisor $p^2$ fails because $p^3>M$. A divisor $pq$ fails in row $q$ because $p^2q>M$. Finally, $pqr$ fails in row $p$ because $(qr)^2>W$, including repeated $Q$-primes.

Therefore
\[
 \K_M^0(b)=\frac1{200}v_p(b)
  \mathbf1_{\{p\mid b,\ \exists q\in Q:q\mid b/p\}}.
\]
For a source $a=pt$, the witness $pq$ is eligible exactly when $q\mid t$ and $q^2\le t$. Such a source has $v_p(a)=1$, since $p\mid t$ would imply $t\ge pq>W$. Consequently
\begin{equation}\label{eq:counter-source}
 S_0:=\sum_{a\le M}\K_a^0(a)
 =\frac1{200}\#\{t\le W:\exists q\in Q, q\mid t, q^2\le t\}.
\end{equation}

\subsection{The source and the exact target mean}
No integer up to $W$ has four distinct $Q$-prime divisors, whereas every triple product is at most $W^{0.96}$. Inclusion-exclusion therefore terminates exactly at degree three. If $e_j$ denotes the elementary symmetric sum of degree $j$ in $\{1/q:q\in Q\}$, the unrestricted union count is
\[
 W(e_1-e_2+e_3)+O(|Q|^3).
\]
The delay $q^2\le t$ removes at most $\sum_{q\in Q}(q-1)=O(W^{0.64})$ integers. Also $|Q|^3=O(W^{0.96})=o(W)$. Mertens' prime harmonic formula gives
\[
 e_1\longrightarrow\mu=\log(32/31),\quad
 e_2\longrightarrow\mu^2/2,\quad e_3\longrightarrow\mu^3/6,
\]
and hence
\begin{equation}\label{eq:counter-source-asym}
 S_0=\frac W{200}\left(\mu-\frac{\mu^2}{2}+\frac{\mu^3}{6}+o(1)\right).
\end{equation}

Choose an actual integer $x\ge0$ with $x\equiv0\pmod{p^2}$ and let
\[
 I_x=\{px+1,\ldots,px+M\}.
\]
Its $p$-multiples are $p(x+t)$ for $1\le t\le W$. Since $W<p^2$,
\[
 v_p\bigl(p(x+t)\bigr)=1+\mathbf1_{p\mid t}
\]
exactly. Average over CRT solutions with independent uniform residues $x\bmod q$, $q\in Q$, keeping $x\equiv0\pmod{p^2}$. For every fixed $t$, the union probability is $1-\prod_{q\in Q}(1-1/q)$, so
\begin{equation}\label{eq:counter-target-exact}
 \mathbb E_x\sum_{b\in I_x}\K_M^0(b)
 =\frac1{200}\left(W+\left\lfloor\frac Wp\right\rfloor\right)
       \left(1-\prod_{q\in Q}(1-1/q)\right).
\end{equation}
No independence between different positions $t$ is needed. Some genuine shift has target sum $T_0$ no greater than this mean. Since the product tends to $e^{-\mu}=31/32$,
\[
 S_0-T_0\ge(\kappa+o(1))W,
\quad
 \kappa=\frac1{200}\left(e^{-\mu}-1+\mu-\frac{\mu^2}{2}+\frac{\mu^3}{6}\right)>0.
\]
Taylor's integral remainder gives
\[
 \frac{e^{-\mu}\mu^4}{4800}\le\kappa\le\frac{\mu^4}{4800},
 \qquad \kappa=2.10334555597958\ldots\times10^{-10}.
\]
The positive constant is very small; no moderate-size numerical witness is claimed.

\subsection{Positive background prices and the empty full tail}
Fix the shift above. Give every remaining prime $r\le M$ the same price $\varepsilon>0$. Every newly available witness contains a background prime, so its price is at most $\varepsilon$. At each \emph{fixed} scale $L$,
\[
 0\le\K_L^\varepsilon(t)-\K_L^0(t)
       \le\varepsilon\Om_{\le M}(t).
\]
The total valuation sum over the finite prefix and the now-fixed target interval is finite. Choose $\varepsilon$ so that the perturbation of either kernel sum is at most one and $\varepsilon\log_2M<\delta/2$. Then $h_\varepsilon(a)=a$ for all sources, and
\[
 \sum_{a\le M}\K_{h_\varepsilon(a)}^\varepsilon(a)
 -\sum_{b\in I_x}\K_M^\varepsilon(b)
 \ge(\kappa+o(1))W-1.
\]
The deficit is also $O(W)$, since $S_0\le W/200$ and the target is nonnegative. Thus it is $\Theta(W)=\Theta(M^{10/17})$, exceeding every fixed multiple of $\sqrt M$.

On the source prefix, $X_0(a)\le2$: $p^3>M$; a source with $p^2$ has no $Q$-factor; a source with one $p$ has at most three $Q$-occurrences; a source without $p$ has at most five. The chosen background gives $X_\varepsilon(a)<2+\delta$. Hence the full tail is empty. All positive source payments in this counterexample are outside the actual full tail. It does not contradict the required estimate on genuine positive optimizers.

\paragraph{A discontinuity exposed by the counterexample.}
The function $w\mapsto\K_{h_w(a)}^w(a)$ is not continuous when a price reaches zero. Infinitesimal positive weights can change $h_w(a)$ and open a main-prime witness of fixed price $1/200$. The fixed-scale estimate above does not justify continuity at the moving active scale. A weighted-demand cleanup must control the raw objective and then rebuild the kernel on the retained support.

\section{A stable corrected quantity at a genuine optimum}\label{sec:pruning}
For the same own full tail $A$ and $n=|A|$, put
\[
 \Phi_A(w)=\sum_{a\in A}X_w(a)-\sum_{b\in I}(X_w(b)-1)_+.
\]
The exact reserve correction is
\[
 \Gamma_c-R_A-R_I=\Phi_A(w)-2n=J_c(w)+\delta n.
\]
Thus the user's sufficient target $\Gamma_c\le R_A+R_I+o_\delta(n)$ is precisely the raw estimate $\Phi_A(w)\le2n+o_\delta(n)$ on this class. Substituting the identity does not prove the estimate. Also this corrected quantity is not $J_2(w)$: the latter includes the additional nonnegative middle-tail sum over $2<X_w(a)\le2+\delta$.

Let $w,z$ be a genuinely paired boxed raw optimum for this fixed $A$, and take any downward change $0\le w'\le w$. Here the boxes belong to the primal variables $z$. The vector $w$ is optimal for the original dual $w\ge0$, with $w\le1$ chosen using the clipping theorem. Thus every positive-price row is tight, including rows with $w_p=1$. Write
\[
 \Delta(t)=X_w(t)-X_{w'}(t),\qquad
 D=\sum_{a\in A}\Delta(a)=\sum_{b\in I}z_b\Delta(b),
\]
where the equality uses the exact whole-prime row balances. Complementary slackness gives the following exact identity:
\begin{equation}\label{eq:pruning}
 \boxed{\Phi_A(w)-\Phi_A(w')
 =\sum_{X_w(b)>1}(1-X_{w'}(b))_+
   +\sum_{X_w(b)=1}z_b\Delta(b).}
\end{equation}
Indeed the old hot hinge decreases by $\min(\Delta(b),X_w(b)-1)$, while old cold and boundary hinges do not decrease. Substitute the equality for $D$, using $z_b=1$ on hot targets and $z_b=0$ on cold targets. In particular,
\[
 0\le\Phi_A(w)-\Phi_A(w')\le D.
\]
This controls the reserve-corrected objective even if individual kernel terms jump. It does not transfer old balances to a new tail or a new optimum. Small individual prices alone do not imply $D=o(n)$.

No arithmetic choice of a downward variation was found that supplies the required uniform full-tail estimate. The correction removes the false unrestricted transfer route; it does not close Erd\H{o}s 708.

\section*{Reference for the existing raw-to-integer transfer}
Ziming Wang, \emph{Raw fractional covers and worst-case leading coefficients}, section of the Erd\H{o}s 708 interval divisor-covers manuscript. Source checked on 24 November 2026:\\
\href{https://github.com/science-prince/erdos-708-interval-divisor-covers/blob/main/papers/main-paper/raw-transfer.tex}{\texttt{raw-transfer.tex} in \texttt{science-prince/erdos-708-interval-divisor-covers}}.
\par\medskip
The prime harmonic asymptotic used in Section~\ref{sec:user-counter} is recorded as formula 27.11.8 in the NIST Digital Library of Mathematical Functions: \url{https://dlmf.nist.gov/27.11.E8}.
\end{document}
